\chapter{Gastrointestinal Disorders}

\section*{Patient focus}
When you have abdominal pain, nausea, vomiting, or bleeding from the digestive tract, imaging helps find the cause. A CT scan is the most commonly used test for acute abdominal problems---it can show blockages, inflammation, infection, and areas where the bowel is not getting enough blood. Different imaging tests are used for different problems: ultrasound is good for the appendix and gallbladder, CT is best for most emergencies, and MRI is useful for Crohn's disease and pelvic conditions.


\label{ch:ch63_gi_disorders}

\section{Imaging}
Radiography, ultrasound, CT, MRI, fluoroscopy, and endoscopy answer different gastrointestinal questions. Select the study according to suspected obstruction, inflammation, bleeding, ischemia, perforation, mass, or postoperative complication.

\section{Bowel Obstruction}

\subsection{Small Bowel Obstruction (SBO)}
\begin{itemize}
  \item \textbf{CT findings}: dilated small bowel ($>$2.5 cm) proximal to a transition point, collapsed distal bowel
  \item \textbf{Closed-loop obstruction}: U-shaped or C-shaped loop with mesenteric swirl---surgical emergency due to strangulation risk
  \item \textbf{Bowel wall thinning}: suggests ischaemia; wall thickening with target sign suggests congestion
  \item \textbf{Common causes}: adhesions (most common in post-surgical patients), hernia, tumour, intussusception, volvulus
\end{itemize}

\subsection{Large Bowel Obstruction (LBO)}
\begin{itemize}
  \item \textbf{CT findings}: dilated colon ($>$6 cm, caecum $>$9 cm) with transition point
  \item \textbf{Most common cause}: colorectal cancer (left side), volvulus (sigmoid or caecal)
  \item \textbf{Caecal perforation risk}: diameter $>$9--12 cm or wall thinning---requires urgent decompression or surgery
\end{itemize}

\begin{pearlbox}[SBO algorithm]
1. CT with IV contrast to confirm SBO and identify transition point.
2. No closed-loop, no ischaemia, partial obstruction $\rightarrow$ trial of conservative management.
3. Closed-loop, complete obstruction, or signs of ischaemia (wall thickening, reduced enhancement, mesenteric oedema, pneumatosis) $\rightarrow$ urgent surgery.
4. Recurrent SBO without clear surgical cause $\rightarrow$ consider adhesiolysis.
\end{pearlbox}

\section{Inflammatory Bowel Disease}

\subsection{Crohn's Disease}
\begin{itemize}
  \item \textbf{MR enterography} (preferred) or CT enterography: distended bowel with neutral oral contrast, IV contrast
  \item \textbf{Key findings}: mural thickening ($>$3 mm), mural oedema (T2 hyperintensity), enhancement, comb sign (engorged vasa recta), skip lesions, mesenteric lymphadenopathy
  \item \textbf{Complications}: fistula, abscess, stricture, phlegmon
  \item \textbf{MR Enterography protocol}: oral contrast to distend bowel, IV gadolinium, T2 fat-sat, DWI, dynamic post-contrast
\end{itemize}

\subsection{Ulcerative Colitis}
\begin{itemize}
  \item \textbf{CT colonography}: continuous mucosal thickening, pericolonic fat stranding, haustral loss (lead-pipe colon in chronic disease)
  \item \textbf{Complications}: toxic megacolon (colonic dilation $>$6 cm with systemic toxicity), perforation, haemorrhage
  \item \textbf{Differential from Crohn's}: UC is continuous, colonic only, mucosal; Crohn's is skip lesions, transmural, any GI segment
\end{itemize}

\begin{table}[h]
\centering
\caption{Crohn's vs Ulcerative Colitis on imaging}
\begin{tabular}{|l|l|l|}
\hline
\textbf{Feature} & \textbf{Crohn's} & \textbf{Ulcerative Colitis} \\
\hline
Comb sign & Present & Absent \\
Creeping fat & Present & Absent \\
Distribution & Skip lesions, any GI segment & Continuous, colon only \\
Fistula/abscess & Common & Rare \\
Lead-pipe colon & Rare & Classic in chronic disease \\
Wall thickening & Focal, eccentric & Diffuse, symmetric \\
\hline
\end{tabular}
\end{table}

\section{Mesenteric Ischaemia}
\begin{itemize}
  \item \textbf{CT with arterial and portal venous phases}: essential for detecting vascular occlusion and bowel changes
  \item \textbf{Vascular findings}: SMA/SMV thrombosis or embolus, non-occlusive mesenteric ischaemia (vascular spasm, low flow state)
  \item \textbf{Bowel findings}: wall thickening (early) or thinning (late), reduced enhancement, pneumatosis intestinalis, portal venous gas, mesenteric oedema, free air
  \item \textbf{Sensitivity}: CT has $>$95\% sensitivity for acute mesenteric ischaemia when bowel changes are present
\end{itemize}

\begin{pearlbox}[When to suspect mesenteric ischaemia]
Think of it in: atrial fibrillation, recent MI, aortic procedure, known mesenteric atherosclerosis, or ``pain out of proportion to examination.'' Early CT may be normal---if clinical suspicion is high, repeat in 6--12 hours or proceed to angiography. Late findings (pneumatosis, portal venous gas) indicate irreversible ischaemia.
\end{pearlbox}

\section{GI Bleeding}

\subsection{Lower GI Bleeding}
\begin{itemize}
  \item \textbf{CT angiography}: active extravasation visible when bleeding rate $>$0.3--0.5 mL/min; localise source
  \item \textbf{Common causes}: diverticular bleeding (most common), angiodysplasia, ischaemic colitis, inflammatory bowel disease, colorectal cancer
  \item \textbf{Nuclear medicine}: Tc-99m labelled RBC scan detects slower bleeding ($>$0.1 mL/min) but has poor localisation
  \item \textbf{Angiography}: therapeutic (embolisation) when CTA shows active bleeding
\end{itemize}

\subsection{Upper GI Bleeding}
\begin{itemize}
  \item \textbf{CT angiography}: rapid assessment, may identify source before endoscopy
  \item \textbf{Endoscopy}: first-line for upper GI bleeding; CT if endoscopy unavailable or fails
\end{itemize}

\section{Celiac Disease}
\begin{itemize}
  \item \textbf{Imaging features}: jejunalisation of ileum (reversal of normal fold pattern), duodenal bulb abnormalities, reduced fold pattern in jejunum
  \item \textbf{CT/MRI complications}: splenic atrophy, enteropathy-associated T-cell lymphoma, small bowel adenocarcinoma, gallstones, associated autoimmune conditions
  \item \textbf{Dental and skeletal}: dental enamel defects, osteoporosis from malabsorption
\end{itemize}

\section{Malrotation and Midgut Volvulus}
\begin{itemize}
  \item \textbf{Paediatric emergency}: most common presentation is bilious vomiting in a neonate
  \item \textbf{Upper GI series}: duodenum crosses midline at D3-D4 (normally fixed to left by Ligament of Treitz); swirling of SMA/SMV on US or CT (whirlpool sign)
  \item \textbf{CT in adults}: may show reversal of SMA/SMV relationship,mesenteric swirling, bowel malposition
  \item \textbf{Treatment}: Ladd's procedure---surgical emergency when volvulus is suspected
\end{itemize}

\section{Interpretation}
Assess bowel caliber and enhancement, transition point, wall, mesentery, vessels, peritoneum, solid organs, and collections. State whether there is obstruction, ischemia, perforation, active bleeding, or a drainable collection.

\section{Safety}
Contrast, pregnancy, aspiration risk, renal function, and procedural risks are individualized. Ischemia, perforation, closed-loop obstruction, severe bleeding, or toxic megacolon requires immediate escalation.

\begin{pearlbox}[GI imaging red flags]
Communicate immediately: closed-loop obstruction, pneumoperitoneum, pneumatosis with portal venous gas, active arterial extravasation, toxic megacolon, and mesenteric vessel occlusion. These findings require urgent surgical or interventional input.
\end{pearlbox}

\clearpage
